% !TeX program = xelatex
% ============================================================
% Linux System Configuration in the AI Era
% XeLaTeX + TikZ
%
% Design goals:
%   1. Keep visual styles and text content separated.
%   2. Use relative node positioning only; no absolute (x,y) coordinates.
%   3. Keep the feedback path away from the responsibility cards.
%   4. Match the companion "AI-assisted C/C++ development workflow".
% ============================================================

\documentclass[10pt,a4paper]{article}

% -------------------- 1. PACKAGES ---------------------------
\usepackage[a4paper,margin=13mm]{geometry}
\usepackage{fontspec}
\usepackage{xeCJK}
\usepackage{xcolor}
\usepackage{tikz}
\usetikzlibrary{
  arrows.meta,
  positioning,
  shapes.geometric,
  fit,
  backgrounds,
  calc
}
\usepackage{enumitem}
\usepackage{microtype}
\pagestyle{empty}

% -------------------- 2. TYPOGRAPHY + STYLE -----------------
% Replace these font lines if another local font is preferred.
\setmainfont{DejaVu Sans}
\setCJKmainfont{Noto Sans CJK SC}
\setsansfont{DejaVu Sans}
\setCJKsansfont{Noto Sans CJK SC}

% ---- Palette: shared with the companion C++ workflow ----
\definecolor{HumanBlue}{HTML}{2563EB}
\definecolor{AIPurple}{HTML}{7C3AED}
\definecolor{ToolGreen}{HTML}{059669}
\definecolor{ReviewOrange}{HTML}{EA580C}
\definecolor{PassGreen}{HTML}{15803D}
\definecolor{FailRed}{HTML}{DC2626}
\definecolor{Ink}{HTML}{1F2937}
\definecolor{SoftGray}{HTML}{F3F4F6}
\definecolor{LineGray}{HTML}{64748B}
\definecolor{PaleBlue}{HTML}{EFF6FF}
\definecolor{PalePurple}{HTML}{F5F3FF}
\definecolor{PaleGreen}{HTML}{ECFDF5}
\definecolor{PaleOrange}{HTML}{FFF7ED}
\definecolor{ApplyTeal}{HTML}{0F766E}
\definecolor{PaleTeal}{HTML}{F0FDFA}

% ---- Global TikZ knobs ----
\tikzset{
  >={Latex[length=2.8mm,width=1.9mm]},
  every node/.style={font=\sffamily},
  workflow/.style={
    draw=Ink!55,
    line width=0.85pt,
    rounded corners=2.5mm,
    text=Ink,
    align=center,
    minimum height=12.5mm,
    text width=72mm,
    inner xsep=5mm,
    inner ysep=2.7mm,
  },
  human/.style={workflow, draw=HumanBlue, fill=PaleBlue},
  ai/.style={workflow, draw=AIPurple, fill=PalePurple},
  tool/.style={workflow, draw=ToolGreen, fill=PaleGreen},
  review/.style={workflow, draw=ReviewOrange, fill=PaleOrange},
  apply/.style={workflow, draw=ApplyTeal, fill=PaleTeal},
  decision/.style={
    diamond,
    aspect=2.35,
    draw=ReviewOrange,
    fill=PaleOrange,
    line width=0.9pt,
    align=center,
    text=Ink,
    inner sep=1.2mm,
    text width=29mm,
  },
  terminal/.style={
    rounded corners=7mm,
    draw=PassGreen,
    fill=PassGreen!8,
    line width=1pt,
    minimum width=48mm,
    minimum height=10mm,
    align=center,
    text=Ink,
  },
  flow/.style={->, draw=LineGray, line width=1.05pt},
  feedback/.style={->, draw=FailRed, line width=1.05pt, dashed},
  passflow/.style={->, draw=PassGreen, line width=1.05pt},
  rolebox/.style={
    rounded corners=2mm,
    draw=Ink!25,
    fill=white,
    line width=0.6pt,
    inner sep=3mm,
    text width=48mm,
    align=left,
    font=\sffamily\small,
  },
  rulebox/.style={
    rounded corners=2mm,
    draw=FailRed!45,
    fill=FailRed!4,
    line width=0.65pt,
    inner sep=3mm,
    text width=48mm,
    align=left,
    font=\sffamily\small,
  },
  phasebg/.style={
    rounded corners=3mm,
    draw=none,
    fill=SoftGray,
    inner sep=4mm,
  },
  tag/.style={
    rounded corners=1.5mm,
    inner xsep=2.2mm,
    inner ysep=1.0mm,
    font=\sffamily\scriptsize\bfseries,
    text=white,
  }
}

% -------------------- 3. CONTENT ----------------------------
% Most later edits should only touch this section.

\newcommand{\DocTitle}{Linux System Configuration in the AI Era}
\newcommand{\DocSubtitle}{人定目标，AI 出变更，工具验状态；分步执行、可测可退，人做最终验收}

\newcommand{\DesiredTitle}{人：定义 Desired State}
\newcommand{\DesiredBody}{明确目标状态、约束、不变量、性能指标、允许的风险与不可破坏条件}

\newcommand{\CurrentTitle}{工具：采集 Current State}
\newcommand{\CurrentBody}{系统版本、启动参数、服务、网络、存储、IRQ、日志与关键运行状态；先取证，后修改}

\newcommand{\PlanTitle}{AI：分析差距并生成 Change Set}
\newcommand{\PlanBody}{基于实际状态给出配置差异、修改文件、命令、影响范围、验证点与候选回滚方案}

\newcommand{\ReviewTitle}{人：Review 变更与回滚边界}
\newcommand{\ReviewBody}{确认依赖关系、故障影响、远程失联风险、重启要求、维护窗口与恢复路径}

\newcommand{\PreflightTitle}{工具：Preflight / Static Validation}
\newcommand{\PreflightBody}{语法检查、配置验证、dry-run、shellcheck / bash -n / systemd-analyze verify 等}

\newcommand{\ApplyTitle}{受控执行：Staged Apply}
\newcommand{\ApplyBody}{备份原配置；小步、幂等、可审计地应用；一次只改变可归因的一组变量}

\newcommand{\ObserveTitle}{工具：采集 Actual State 与证据}
\newcommand{\ObserveBody}{systemctl / journalctl / dmesg / procfs / sysfs / perf / cyclictest 等重新测量实际状态}

\newcommand{\CompareTitle}{人 + AI：Desired vs. Actual}
\newcommand{\CompareBody}{AI 辅助解释日志和差异；人判断目标是否真正达到，以及副作用是否可接受}

\newcommand{\DecisionText}{目标、证据与\\风险均满足？}
\newcommand{\PassText}{通过：固化配置 / 版本化 / 文档化}
\newcommand{\FailText}{把偏差、日志和测量证据\\反馈给 AI，重新生成变更计划}

\newcommand{\HumanRole}{%
  \textbf{人负责：}\\[-1mm]
  \begin{itemize}[leftmargin=*,nosep]
    \item 定义 Desired State / Invariant
    \item 判断风险、依赖与维护窗口
    \item 决定是否执行、是否回滚
    \item 最终验收与上线责任
  \end{itemize}}

\newcommand{\AIRole}{%
  \textbf{AI 负责：}\\[-1mm]
  \begin{itemize}[leftmargin=*,nosep]
    \item 分析 Current / Actual State
    \item 生成 Change Set 与脚本草案
    \item 辅助解释日志、差异与异常
    \item 生成验证项、文档与回滚建议
  \end{itemize}}

\newcommand{\ToolRole}{%
  \textbf{工具负责：}\\[-1mm]
  \begin{itemize}[leftmargin=*,nosep]
    \item 提供真实系统状态和可复现证据
    \item 做语法、静态和运行期验证
    \item 测量性能、服务状态和副作用
  \end{itemize}}

\newcommand{\SafetyRule}{%
  \textbf{核心规则：}\\[-1mm]
  不让 AI 从“问题描述”直接跳到 root 命令。\\
  至少经过：\textbf{Current State → Change Plan → Verify Plan}。}

% -------------------- 4. DRAWING ----------------------------
\begin{document}

\begin{center}
  {\LARGE\bfseries\color{Ink}\DocTitle}\par
  \vspace{1.2mm}
  {\small\color{LineGray}\DocSubtitle}
\end{center}

\vspace{1.5mm}

\begin{tikzpicture}[node distance=5.1mm and 10mm]

  % ---- Main vertical flow: all positions are relative ----
  \node[human] (desired) {\textbf{\DesiredTitle}\\[0.8mm]\small \DesiredBody};

  \node[tool, below=of desired] (current) {\textbf{\CurrentTitle}\\[0.8mm]\small \CurrentBody};

  \node[ai, below=of current] (plan) {\textbf{\PlanTitle}\\[0.8mm]\small \PlanBody};

  \node[review, below=of plan] (review) {\textbf{\ReviewTitle}\\[0.8mm]\small \ReviewBody};

  \node[tool, below=of review] (preflight) {\textbf{\PreflightTitle}\\[0.8mm]\small \PreflightBody};

  \node[apply, below=of preflight] (apply) {\textbf{\ApplyTitle}\\[0.8mm]\small \ApplyBody};

  \node[tool, below=of apply] (observe) {\textbf{\ObserveTitle}\\[0.8mm]\small \ObserveBody};

  \node[human, below=of observe] (compare) {\textbf{\CompareTitle}\\[0.8mm]\small \CompareBody};

  \node[decision, below=7mm of compare] (decision) {\textbf{\DecisionText}};

  \node[terminal, below=7mm of decision] (pass) {\textbf{\PassText}};

  % ---- Main flow arrows ----
  \draw[flow] (desired) -- (current);
  \draw[flow] (current) -- (plan);
  \draw[flow] (plan) -- (review);
  \draw[flow] (review) -- (preflight);
  \draw[flow] (preflight) -- (apply);
  \draw[flow] (apply) -- (observe);
  \draw[flow] (observe) -- (compare);
  \draw[flow] (compare) -- (decision);
  \draw[passflow] (decision) -- node[right, font=\scriptsize\bfseries, text=PassGreen] {是} (pass);

  % ---- Failure card and return path ----
  % The card is to the right of the decision. The dashed return path travels
  % below the main flow and then up the LEFT side to the AI planning node,
  % so it never crosses the responsibility cards on the right.
  \node[review, text width=49mm, right=18mm of decision] (feedback) {\textbf{发现偏差 / 副作用}\\[0.8mm]\small \FailText};
  \draw[feedback] (decision.east) -- node[above, font=\scriptsize\bfseries, text=FailRed] {否} (feedback.west);

  \coordinate (feedbackbottom) at ([yshift=-7mm]pass.south);
  \coordinate (feedbackleft)   at ([xshift=-9mm]plan.west);
  \draw[feedback]
    (feedback.south)
    |- (feedbackbottom)
    -| (feedbackleft)
    -- node[above, near end, font=\scriptsize\bfseries, text=FailRed] {证据驱动再规划} (plan.west);

  % ---- Responsibility cards ----
  \node[rolebox, right=18mm of current] (humanrole) {\HumanRole};
  \node[rolebox, below=4.5mm of humanrole] (airole) {\AIRole};
  \node[rolebox, below=4.5mm of airole] (toolrole) {\ToolRole};
  \node[rulebox, below=4.5mm of toolrole] (rule) {\SafetyRule};

  \node[tag, fill=HumanBlue, above left=-1mm and 0mm of humanrole.north west] {HUMAN};
  \node[tag, fill=AIPurple, above left=-1mm and 0mm of airole.north west] {AI};
  \node[tag, fill=ToolGreen, above left=-1mm and 0mm of toolrole.north west] {TOOLS};
  \node[tag, fill=FailRed, above left=-1mm and 0mm of rule.north west] {RULE};

  % ---- Soft background grouping ----
  \begin{scope}[on background layer]
    \node[phasebg, fit=(desired)(current)] {};
    \node[phasebg, fit=(plan)(review)] {};
    \node[phasebg, fit=(preflight)(apply)(observe)] {};
    \node[phasebg, fit=(compare)(decision)(pass)(feedback)] {};
  \end{scope}

\end{tikzpicture}

\vfill
\begin{center}
  \small\bfseries\color{Ink}
  Desired State → Evidence → Change Set → Verify → Staged Apply → Observe → Review
\end{center}

\end{document}
